\documentclass[11pt,a4paper]{article}
\usepackage[T1]{fontenc}
\usepackage[utf8]{inputenc}
\usepackage[margin=27mm,headheight=15pt]{geometry}
\usepackage{amsmath,amsthm,mathtools}
\usepackage{newtxtext,newtxmath}
\usepackage{microtype}
\usepackage{booktabs,array,longtable,tabularx}
\usepackage{xcolor}
\usepackage{enumitem}
\usepackage{listings}
\usepackage[most]{tcolorbox}
\usepackage{fancyhdr}
\usepackage{hyperref}
\usepackage{bookmark}
\definecolor{ink}{HTML}{354139}
\definecolor{quiet}{HTML}{626B5C}
\definecolor{paperbox}{HTML}{F4F5EF}
\hypersetup{colorlinks=true,linkcolor=ink,citecolor=ink,urlcolor=ink,
 pdftitle={Real-Branch Asymptotic Reversion with Logarithms and Irrational Powers},
 pdfauthor={Prepared for Vladimir Reshetnikov},pdfsubject={Theory and Wolfram Language implementation}}
\pagestyle{fancy}\fancyhf{}
\fancyhead[L]{\small\scshape Real-branch asymptotic reversion}
\fancyhead[R]{\small PowerLogInverse 1.0}
\fancyfoot[C]{\thepage}
\setlength{\parskip}{0.35em}
\setlength{\parindent}{1.1em}
\setcounter{tocdepth}{2}
\numberwithin{equation}{section}
\allowdisplaybreaks[2]
\newtheorem{theorem}{Theorem}[section]
\newtheorem{proposition}[theorem]{Proposition}
\newtheorem{lemma}[theorem]{Lemma}
\newtheorem{corollary}[theorem]{Corollary}
\theoremstyle{definition}
\newtheorem{definition}[theorem]{Definition}
\newtheorem{example}[theorem]{Example}
\theoremstyle{remark}
\newtheorem{remark}[theorem]{Remark}
\newcommand{\N}{\mathbb N}
\newcommand{\R}{\mathbb R}
\newcommand{\C}{\mathbb C}
\newcommand{\dd}{\mathrm d}
\newcommand{\OO}{\mathcal O}
\newcommand{\val}{\operatorname{val}}
\newcommand{\supp}{\operatorname{supp}}
\newcommand{\coeff}[2]{\left[#1\right]#2}
\newcommand{\pl}{\texttt{PowerLogInverse}}
\newcommand{\code}[1]{\texttt{#1}}
\newtcolorbox{resultbox}[1][]{enhanced,breakable,colback=paperbox,colframe=quiet,
 boxrule=0.45pt,arc=0pt,left=8pt,right=8pt,top=6pt,bottom=6pt,#1}
\lstdefinelanguage{WL}{morekeywords={Module,Table,Total,Do,If,Return,FullSimplify,
 D,Log,Exp,ProductLog,PowerLogInverse,PowerLogInverseData,PowerLogReversion,
 PowerLogResidual,PowerLogOrder,LogPowerBaseInverse,InversePerturbationJet,Get,
 TestReport,FileNameJoin,DirectoryName},sensitive=true,
 morecomment=[s]{(*}{*)},morestring=[b]"}
\lstset{language=WL,basicstyle=\ttfamily\small,keywordstyle=\color{ink}\bfseries,
 commentstyle=\color{quiet},stringstyle=\color{quiet},backgroundcolor=\color{paperbox},
 frame=single,rulecolor=\color{quiet!40},framerule=0.3pt,breaklines=true,
 columns=fullflexible,keepspaces=true,showstringspaces=false,tabsize=2,
 aboveskip=0.7em,belowskip=0.7em}
\begin{document}
\begin{titlepage}
\vspace*{14mm}
{\color{quiet}\large MATHEMATICAL THEORY AND IMPLEMENTATION\par}
\vspace{8mm}
{\Huge\bfseries Real-Branch Asymptotic Reversion\par}
\vspace{4mm}
{\LARGE with Logarithms and Irrational Powers\par}
\vspace{9mm}
{\large A constructive coefficient calculus, rigorous remainder bounds,\par
 and the Wolfram Language package \texttt{PowerLogInverse}\par}
\vspace{12mm}
{\large Prepared for Vladimir Reshetnikov\par}
{\large 7 September 2026\par}
\vspace{13mm}
\begin{resultbox}
For
\[
 f(x)=a x^p\left(1+\sum_{i=1}^{s}x^{\delta_i}P_i(\log x)\right),
 \qquad a,p,\delta_i>0,
\]
put $u=(y/a)^{1/p}$ and $L=\log u$. The positive inverse tending to zero is
\[
 f^{-1}(y)=u\left(1+\sum_{\boldsymbol m\ne0}
 u^{\boldsymbol m\cdot\boldsymbol\delta}Q_{\boldsymbol m}(L)\right),
\]
where, with $n=|\boldsymbol m|$ and $\Delta=\boldsymbol m\cdot\boldsymbol\delta$,
\[
 Q_{\boldsymbol m}(L)=
 \frac{(-1)^n}{p^n\prod_i m_i!}
 \prod_{j=1}^{n-1}\left(\frac{\dd}{\dd L}+1+\Delta+pj\right)
 \prod_i P_i(L)^{m_i}.
\]
Every finite power cutoff is effectively computable. The logarithmic
factors in its remainder are tracked rather than hidden in an ordinary
power-series order symbol.
\end{resultbox}
\vfill
{\small Package version 1.0.0. The archive includes exact coefficient code,
 an independent sparse residual checker, examples, native Wolfram test
 definitions, and executed Python validation. Native Wolfram-kernel execution
 was not available in the build environment.\par}
\end{titlepage}

\begin{abstract}
The inverse of $x+x^2(1+\log x)$ on the positive real axis has logarithmic
coefficient polynomials, whereas the inverse of $x+x^{\sqrt2}$ has an
irrational exponent progression. Neither example should be treated as an
unspecified analytic inverse at an arbitrary zero of a complexified
function. We explain the complex Taylor germ recorded in the motivating
Mathematica question, then construct the desired real germs directly.
A specialization of the classical Lagrange--B\"urmann formula gives explicit
multi-index differential polynomials for finite sums of power--logarithmic
corrections to $a x^p$. A finitely generated positive exponent semigroup
makes strict cutoff computation finite even when the corresponding
additive group is dense. An analytic implicit-function argument proves
actual convergence for exact finite input of this class, and also gives a
computable asymptotic remainder bound for every cutoff. Residuals transfer
to inverse errors through an explicit derivative estimate. The two original
examples are developed to arbitrary order; the irrational-power example
also has a closed coefficient formula and an exact convergence radius.
Extensions include observable inversion, leading logarithmic pivots with
real Lambert-$W$ branches, uncertain input jets, and flat perturbations.
The accompanying package implements the finite power--logarithmic core,
exact algebraic support ordering, branch-aware normalization, and an
independent truncated residual calculation. Mathematical validation and
native-kernel verification are reported separately.
\end{abstract}
{\setlength{\parskip}{0pt}\small\tableofcontents}
\newpage

\section{The question is about a germ, not merely an inverse symbol}
\label{sec:question}

The archived question~\cite{question} asks for the inverse near $y=0^+$ of
\begin{equation}\label{eq:motivating}
 f_1(x)=x+x^2(1+\log x),\qquad x>0,
\end{equation}
and, in its update, for the corresponding inverse of
\begin{equation}\label{eq:irrinput}
 f_2(x)=x+x^{\sqrt2},\qquad x>0.
\end{equation}
We use $x$ for the source variable and $y$ for the argument of the inverse.
This avoids confusing a series for $f$ with a series for $f^{-1}$.
All logarithms of positive quantities below are real logarithms.

The request has three logically distinct components: identify the intended
inverse branch, identify a scale closed under the necessary operations,
and compute coefficients with a meaningful remainder. Merely attaching
$y>0$ to an unspecified inverse expression addresses none of these
components completely. An inverse at a complex zero can be evaluated for
small positive $y$ and still remain complex.

\subsection{The real branches exist globally in both examples}

For the first function,
\[
 f_1'(x)=1+x(3+2\log x).
\]
The minimum of $x(3+2\log x)$ occurs at $x=e^{-5/2}$ and equals
$-2e^{-5/2}$. Consequently,
\begin{equation}\label{eq:global-slope}
 f_1'(x)\ge 1-2e^{-5/2}>0\qquad(x>0).
\end{equation}
Together with $f_1(x)\to0$ as $x\to0^+$ and $f_1(x)\to\infty$ as
$x\to\infty$, this proves that $f_1$ is a bijection of $(0,\infty)$ onto
itself. For $f_2$, the same conclusion follows from
$f_2'(x)=1+\sqrt2\,x^{\sqrt2-1}>0$ and the endpoint limits.
Thus the real inverse is unambiguous in these examples. The problem is
not a lack of a real branch; it is how to represent and select it.

\subsection{Why the archived output is a correct complex Taylor series}

The archived computation returned a constant involving
\[
 w=W_{-1}(-e),\qquad z_0=e^{w-1}.
\]
Since $-e<-1/e$, this value of $W_{-1}$ is not real. On the compatible local
logarithm branch, $\log z_0=w-1$. Also $we^w=-e$ gives $z_0w=-1$, hence
\[
 f_1(z_0)=z_0\bigl(1+z_0(1+\log z_0)\bigr)=0.
\]
The derivative identities simplify to
\[
 f_1'(z_0)=z_0-1,\qquad f_1''(z_0)=3+2w.
\]
The ordinary analytic inverse at this nonzero complex zero therefore has
\begin{equation}\label{eq:complex-germ}
 G(y)=z_0+\frac{y}{z_0-1}
 -\frac{3+2w}{2(z_0-1)^3}y^2+\OO(y^3),
\end{equation}
exactly explaining the structure of the recorded output. This is a valid
local inverse calculation at the wrong germ for the requested application.
Other Lambert-$W$ indices require attention to the logarithm sheet; one
must not indiscriminately treat every such index as a zero for a fixed
principal logarithm.

Our construction selects a different germ from the outset:
\begin{equation}\label{eq:branch-address}
 g(y)>0,\qquad g(y)\longrightarrow0,\qquad \frac{g(y)}{y}\longrightarrow1.
\end{equation}
For the general leading term $a x^p$, the last condition becomes
$g(y)/(y/a)^{1/p}\to1$. This asymptotic address is part of the mathematical
problem, not a formatting assumption appended to an already chosen branch.

The archive is evidence about a particular recorded computation, not a
benchmark of every current Wolfram Language version. The package avoids
reliance on the branch-selection behavior of \code{InverseFunction}.
The official documentation describes that symbol and its treatment of
inverse branches~\cite{wlinverse}; no new claim about a current built-in
failure is needed here.

\subsection{The answers in their natural scales}

Writing $L=\log y$, the first inverse begins
\begin{align}\label{eq:loganswer}
 g_1(y)={}&y-(L+1)y^2+(2L^2+5L+3)y^3\notag\\
 &-\left(5L^3+\frac{41}{2}L^2+27L+\frac{23}{2}\right)y^4
 +\OO\!\left(y^5(1+|L|)^4\right).
\end{align}
In particular, the tempting two-term notation
$y-(1+\log y)y^2+\OO(y^3)$ is \emph{not} a valid remainder for this real
inverse. The correct two-term estimate is
\begin{equation}\label{eq:twotermerror}
 g_1(y)=y-(1+\log y)y^2+\OO(y^3\log^2 y).
\end{equation}
This observation concerns the inverse remainder; an order symbol used to
truncate the \emph{input} function is a separate matter.

For $r=\sqrt2$, the second inverse begins
\begin{align}\label{eq:irranswer}
 g_2(y)={}&y-y^r+r y^{2r-1}
 -\frac{r(3r-1)}2y^{3r-2}
 +\OO(y^{4r-3})\notag\\
 ={}&y-y^{\sqrt2}+\sqrt2\,y^{2\sqrt2-1}
 -\frac{6-\sqrt2}{2}y^{3\sqrt2-2}
 +\OO(y^{4\sqrt2-3}),
\end{align}
which is the four-block expansion requested in the update.

The first inverse has a $C^1$ extension to zero but not a $C^2$ one:
$g_1''(y)\sim-2\log y\to+\infty$. For the second,
$g_2''(y)\sim-r(r-1)y^{r-2}\to-\infty$. Ordinary Taylor expansions at zero
cannot encode these behaviors. Ordinary Puiseux series on a fixed rational
power grid do not encode the irrational progression either.

\section{Power--log blocks and cutoff-finite support}
\label{sec:support}

\subsection{The representation}

A power--log block is an expression
\[
 u^\gamma Q(L),\qquad L=\log u,\quad Q\in\R[L].
\]
We keep a whole logarithmic polynomial at a given power together. For
fixed $k,\ell$ and $\gamma<\eta$,
\[
 u^\eta(\log u)^\ell=o\bigl(u^\gamma(\log u)^k\bigr)
 \qquad(u\to0^+),
\]
provided the denominator monomial is nonzero. Thus the exponent is the
primary order coordinate; within a block the largest nonzero logarithmic
degree determines its leading magnitude. Odd logarithmic powers can be
negative, so magnitude estimates use $|\log u|$.

Let $\delta_1,\ldots,\delta_s$ be fixed positive real numbers and define
\begin{equation}\label{eq:semigroup}
 \Gamma=\left\{\sum_{i=1}^s m_i\delta_i:m_i\in\N\right\}.
\end{equation}
The formal relative corrections live in locally finite sums
$\sum_{\gamma\in\Gamma}u^\gamma Q_\gamma(L)$.
For the inverse itself the support lies in $1+\Gamma$.

\begin{lemma}[Local finiteness]\label{lem:finite}
For every finite $B$, only finitely many multi-indices satisfy
$\boldsymbol m\cdot\boldsymbol\delta<B$. In particular, $\Gamma\cap[0,B)$
is finite.
\end{lemma}
\begin{proof}
Every such index satisfies $0\le m_i<B/\delta_i$. It is therefore contained
in a finite integer box. Different indices may have the same weight,
which only reduces the number of distinct support points.
\end{proof}

This distinction is important. The additive \emph{group} generated by two
incommensurable real numbers can be dense, but our \emph{positive semigroup}
is locally finite. Irrationality does not create infinitely many terms
below a fixed cutoff. It can create successively smaller gaps far out in
the support; no global lower bound on those gaps is assumed.

\subsection{Formal arithmetic and composition}

For a nonzero formal sum, define its power valuation to be its least
exponent. Multiplication adds valuations. A logarithmic polynomial is a
coefficient and has power valuation zero. Differentiation obeys
\begin{equation}\label{eq:Euler}
 \frac{\dd}{\dd u}\bigl(u^\gamma Q(\log u)\bigr)
 =u^{\gamma-1}(\gamma+D_L)Q(L),\qquad D_L=\frac{\dd}{\dd L}.
\end{equation}
One may enlarge the ambient exponent group when differentiating, but
inside every algorithm below only a controlled translate of $\Gamma$
is used.

If $V$ has positive valuation, the formal expansions
\begin{align}\label{eq:unit-calculus}
 (1+V)^q&=\sum_{k\ge0}\binom qk V^k,\notag\\
 \log(1+V)&=\sum_{k\ge1}\frac{(-1)^{k+1}}k V^k
\end{align}
are cutoff-finite: only finitely many powers $V^k$ can affect a prescribed
cutoff. The substitution $x=u(1+V)$ is consequently well defined by
\[
 x^\delta=u^\delta(1+V)^\delta,\qquad
 \log x=L+\log(1+V).
\]
These are identities on the positive real branch when the formal sums
are realized by sufficiently small actual corrections.

\subsection{Why not manufacture a \texttt{SeriesData} object?}

The documented \code{SeriesData} representation indexes powers by
$n/\mathrm{den}$~\cite{wlseries}. A custom sparse list of exact exponents
and logarithmic coefficient polynomials makes the intended support and
cutoff explicit without relying on that representation to perform
nonstandard reversion. The package therefore returns ordinary finite
expressions together with structured metadata. Its remainder object is
an inert descriptor, not a counterfeit \code{SeriesData} object.

This finite-support viewpoint is compatible with broader transseries
frameworks~\cite{vdh}. The supplied ProveIt draft corpus~\cite{proveit}
also emphasizes support-aware operations, real-germ selection, and the
separation between formal calculation and analytic error transfer. Those
principles are useful here without importing an entire general transseries
field into the implementation.

\section{Prepared finite germs and the selected real inverse}
\label{sec:prepared}

\begin{definition}[Prepared finite power--log germ]\label{def:class}
A prepared finite germ is
\begin{equation}\label{eq:class}
 f(x)=a x^p(1+H(x)),\qquad
 H(x)=\sum_{i=1}^s x^{\delta_i}P_i(\log x),
\end{equation}
where $a,p>0$, every $\delta_i>0$, and every $P_i$ is a real polynomial.
Equal exponents may be combined and zero polynomials removed. The case
$s=0$ is allowed and is an exact monomial.
\end{definition}

The theory permits arbitrary positive real exponents. The executable
package restricts support exponents to exact real algebraic numbers so
that order and equality decisions are exact and operationally meaningful.

Since a positive power beats every fixed logarithmic power,
\[
 H(x)\to0,\qquad xH'(x)\to0.
\]
Consequently,
\begin{equation}\label{eq:derivativeasym}
 f'(x)=a x^{p-1}\bigl(p(1+H(x))+xH'(x)\bigr)
 \sim apx^{p-1}.
\end{equation}

\begin{proposition}[Existence and branch address]\label{prop:branch}
Every prepared finite germ has a unique positive real inverse on some
interval $(0,y_0)$ with $g(y)\to0$. With
\begin{equation}\label{eq:uniformizer}
 u=\left(\frac ya\right)^{1/p}>0,\qquad
 L=\log u=\frac1p\log\frac ya,
\end{equation}
one has $g(y)/u\to1$.
\end{proposition}
\begin{proof}
Equation~\eqref{eq:derivativeasym} makes $f$ strictly increasing and positive
on a sufficiently short interval $(0,x_0)$. Its limit at zero is zero,
so it maps this interval bijectively onto $(0,f(x_0))$. If $x=g(y)$, then
$u=x(1+H(x))^{1/p}$, whence $x/u\to1$.
\end{proof}

This is an eventual local assertion. Unlike the two motivating examples,
a general prepared finite germ need not be globally monotone or injective.
For example, a large negative correction can create turning points away
from zero. The package returns the addressed small positive branch and
does not claim to classify remote branches.

\begin{remark}[Cutoff convention]
A package cutoff $C>1$ means: retain inverse blocks whose \emph{power of
$u$} is strictly less than $C$. In the original target variable this is a
strict $y$-power cutoff $C/p$. Thus for $p\ne1$ a cutoff number must not be
read as an exponent of $y$. The normalization rules are included in the
returned data specifically to expose this distinction.
\end{remark}

\section{An explicit multi-index reversion formula}
\label{sec:formula}

For $\boldsymbol m=(m_1,\ldots,m_s)\in\N^s$, put
\[
 n=|\boldsymbol m|=\sum_i m_i,\qquad
 \Delta=\sum_i m_i\delta_i,\qquad
 \boldsymbol m!=\prod_i m_i!,\qquad
 P^{\boldsymbol m}(L)=\prod_i P_i(L)^{m_i}.
\]
An empty product of differential operators is the identity.

\begin{theorem}[Power--logarithmic inverse coefficients]\label{thm:main}
The addressed inverse of~\eqref{eq:class} has the expansion
\begin{equation}\label{eq:mainseries}
 g(y)=u\left(1+\sum_{\boldsymbol m\in\N^s\setminus\{0\}}
 u^\Delta Q_{\boldsymbol m}(L)\right),
\end{equation}
where
\begin{equation}\label{eq:Qformula}
 \boxed{\displaystyle
 Q_{\boldsymbol m}(L)=
 \frac{(-1)^n}{p^n\boldsymbol m!}
 \left[\prod_{j=1}^{n-1}(D_L+1+\Delta+pj)\right]
 P^{\boldsymbol m}(L).}
\end{equation}
For exact finite input this series converges absolutely for all sufficiently
small positive $u$. As a formal series it is the unique inverse with
leading term $u$. Terms from different multi-indices with the same $\Delta$
must be added before an actual nonzero support is reported.
\end{theorem}

The coefficient identity is a specialization of classical
Lagrange--B\"urmann inversion, not a claim of a new general inversion
theorem. Its useful features here are the explicit logarithmic differential
operators, the real normalization, and a cutoff mechanism adapted to
irrational support.

\subsection{Homotopy form of Lagrange--B\"urmann}

Let $Z+\tau A(Z)=Y$, and let $\Phi$ be analytic near a nonzero base value
$Y$. The inverse germ satisfying $Z=Y$ at $\tau=0$ obeys
\begin{equation}\label{eq:LB}
 \Phi(Z)=\Phi(Y)+\sum_{n\ge1}\frac{(-\tau)^n}{n!}
 \frac{\dd^{n-1}}{\dd Y^{n-1}}
 \bigl(\Phi'(Y)A(Y)^n\bigr).
\end{equation}
This is the usual observable form of the Lagrange inversion theorem;
standard formulations and residue versions are given in
DLMF~\S1.10(vii)~\cite{dlmfinverse}. Importantly, the expansion variable
in~\eqref{eq:LB} is initially the homotopy parameter $\tau$, not $Y$ at a
singular origin. For fixed positive $Y$, real powers and logarithms are
analytic on a small complex neighborhood with the chosen branches.
The uniform justification for setting $\tau=1$ as $Y\to0^+$ will be proved
in Section~\ref{sec:convergence}.

\subsection{Reduction to a near-identity map}

Set
\[
 Z=x^p,\qquad Y=y/a=u^p,\qquad \Phi(Z)=Z^{1/p}.
\]
Then $y=f(x)$ is equivalent to
\[
 Y=Z+A(Z),\qquad
 A(Z)=Z\sum_i Z^{\delta_i/p}
 P_i\!\left(\frac{\log Z}{p}\right).
\]
Expanding $A^n$ by the multinomial theorem gives, for an index of total
size $n$,
\[
 \Phi'(Y)\,\frac{n!}{\boldsymbol m!}
 Y^{n+\Delta/p}P^{\boldsymbol m}(L)
 =\frac{n!}{p\boldsymbol m!}
 Y^{n-1+(1+\Delta)/p}P^{\boldsymbol m}(L).
\]
The elementary derivative identity
\begin{equation}\label{eq:derivY}
 \frac{\dd^k}{\dd Y^k}\left(Y^b Q\!\left(\frac{\log Y}{p}\right)\right)
 =Y^{b-k}\prod_{h=0}^{k-1}\left(b-h+\frac1pD_L\right)Q(L)
\end{equation}
follows by induction. Apply it with
$k=n-1$ and $b=n-1+(1+\Delta)/p$. Reversing the order of the commuting
factors and using $Y^{(1+\Delta)/p}=u^{1+\Delta}$ yields exactly
\eqref{eq:Qformula}.

The operators commute because they are all $D_L$ plus scalar constants.
They preserve polynomiality and do not increase logarithmic degree. Thus
\begin{equation}\label{eq:degreebound}
 \deg Q_{\boldsymbol m}\le\sum_i m_i\deg P_i.
\end{equation}
Cancellation can lower the degree or annihilate a block completely.

\subsection{The first two correction levels}

For $\boldsymbol m=\boldsymbol e_i$,
\[
 Q_{\boldsymbol e_i}=-\frac{P_i}{p}.
\]
For a repeated correction,
\[
 Q_{2\boldsymbol e_i}
 =\frac{(D_L+1+2\delta_i+p)P_i^2}{2p^2}.
\]
For $i\ne j$,
\[
 Q_{\boldsymbol e_i+\boldsymbol e_j}
 =\frac{(D_L+1+\delta_i+\delta_j+p)(P_iP_j)}{p^2}.
\]
The factor $1/2$ appears only in the repeated case. It comes from
$\boldsymbol m!$, not from a blanket $n!$ after multinomial expansion.
This is an easy place for an implementation to introduce a mixed-term bug.

\subsection{Formal uniqueness and a triangular alternative}

Write $g=u(1+V)$ and substitute into $f(g)=a u^p$. The relative equation is
\begin{equation}\label{eq:Veq}
 \mathcal F(V)=
 (1+V)^p\left[1+\sum_i u^{\delta_i}(1+V)^{\delta_i}
 P_i\bigl(L+\log(1+V)\bigr)\right]-1=0.
\end{equation}
Its linear term at power valuation zero is $pV$. All additional terms
involving $V$ either contain a positive support weight or are at least
quadratic in $V$. If two solutions first differ at weight $\gamma$, the
coefficient of their difference in~\eqref{eq:Veq} is $p$ times that first
block. Since $p\ne0$, the block must vanish, a contradiction.

This also gives a recursive algorithm: visit support weights in increasing
order; at weight $\gamma$, evaluate the already determined residual and
set the new coefficient equal to minus that residual coefficient divided
by $p$. The closed formula and this triangular procedure solve the same
formal problem by different computational routes.

\section{Convergence and explicit asymptotic remainders}
\label{sec:convergence}

A formal identity does not by itself supply a real asymptotic theorem.
Here the finite nature of the input makes an analytic proof particularly
clean.

\subsection{Desingularizing the logarithmic coefficients}

Write $P_i(L)=\sum_{j=0}^{d_i}c_{ij}L^j$. Introduce independent small
parameters
\[
 \zeta_{ik}=u^{\delta_i}L^k,\qquad 0\le k\le d_i.
\]
In the normalized equation for $v=g/u$, replace the actual combinations
$u^{\delta_i}L^k$ by these independent parameters. Expanding
$P_i(L+\log v)$ gives the analytic equation
\begin{equation}\label{eq:implicitfinite}
 v^p\left[
 1+\sum_i\sum_{j=0}^{d_i}\sum_{k=0}^{j}
 c_{ij}\binom jk\zeta_{ik}v^{\delta_i}(\log v)^{j-k}
 \right]=1.
\end{equation}
Near $v=1$, fix the analytic logarithm with $\log1=0$; all real powers of
$v$ are then analytic. At $\boldsymbol\zeta=0$, the solution is $v=1$
and the derivative with respect to $v$ is $p\ne0$.

The analytic implicit-function theorem therefore gives a convergent
multivariate Taylor expansion $v=V(\boldsymbol\zeta)$ in a fixed
polydisc around the origin. Every actual parameter $u^{\delta_i}L^k$
tends to zero as $u\to0^+$, so the real curve of parameters eventually
lies in that polydisc. The resulting solution is the addressed inverse
from Proposition~\ref{prop:branch}.

Grouping the absolutely convergent multivariate Taylor series by the
number of factors originating from each correction index $i$ gives
\eqref{eq:mainseries}. For each fixed multi-index, only finitely many
logarithmic choices occur. The coefficient computation in
Section~\ref{sec:formula} identifies the resulting polynomial with
\eqref{eq:Qformula}. This proves the convergence assertion in
Theorem~\ref{thm:main}, not merely formal consistency.

\subsection{A computable remainder for every strict cutoff}

For $C>1$, put $B=C-1$ and define the finite approximant
\begin{equation}\label{eq:AC}
 A_C(y)=u\left(1+\sum_{0<\Delta<B}
 u^\Delta Q_{\boldsymbol m}(L)\right),
\end{equation}
where the summation is over multi-indices and equal powers are collected.
Let
\begin{equation}\label{eq:Nandd}
 \delta_* =\min_i\delta_i,\qquad
 d=\max_i\deg P_i,\qquad N=\left\lceil\frac{B}{\delta_*}\right\rceil.
\end{equation}

\begin{theorem}[Uniformizer-power cutoff bound]\label{thm:remainder}
For nonempty exact finite input,
\begin{equation}\label{eq:remainder}
 g(y)-A_C(y)=\OO\!\left(u^C(1+|\log u|)^{Nd}\right)
 \qquad(y\to0^+).
\end{equation}
For an exact monomial, the inverse is exactly $u$ and the remainder is zero.
\end{theorem}
\begin{proof}
Truncate the analytic multivariate Taylor series
$V(\boldsymbol\zeta)$ at total degree $N-1$. Its remaining tail is
$\OO(\|\boldsymbol\zeta\|^N)$ in a smaller fixed polydisc. Along the
actual curve,
\[
 \|\boldsymbol\zeta\|
 =\OO\!\left(u^{\delta_*}(1+|L|)^d\right),
\]
so this tail is
$\OO(u^{N\delta_*}(1+|L|)^{Nd})$. Since $N\delta_*\ge B$, it is
$\OO(u^B(1+|L|)^{Nd})$.

Among the finitely many terms of total degree below $N$, those omitted
from~\eqref{eq:AC} have weight $\Delta\ge B$ and logarithmic degree at
most $(N-1)d$. Their sum obeys the same bound. No term of total degree
$N$ or larger can have weight below $B$, because
$\Delta\ge N\delta_*\ge B$. Multiplying the relative estimate by $u$
proves~\eqref{eq:remainder}.
\end{proof}

The bound is intentionally conservative for mixed generators. It requires
no search for the first nonzero omitted shell and survives arbitrary
coefficient cancellations. The package returns precisely this bound.
In the one-generator logarithmic example with integer cutoff $C$, it has
the sharp logarithmic degree $C-1$ for the first omitted block.

\begin{corollary}[Pure-power weakening]
For every fixed $\varepsilon>0$,
\[
 g(y)-A_C(y)=\OO(u^{C-\varepsilon}).
\]
One cannot generally replace this by $\OO(u^C)$ when logarithmic
coefficients occur at the boundary.
\end{corollary}

\subsection{Sharper next-shell estimates}

Suppose $\lambda\ge B$ is the least omitted weight with a nonzero
\emph{collected} polynomial $Q_\lambda$. Local finiteness provides a gap
to the next possible weight. Applying Theorem~\ref{thm:remainder} at a
cutoff inside that gap gives
\[
 g(y)-A_C(y)
 =u^{1+\lambda}Q_\lambda(L)
 +o\!\left(u^{1+\lambda}(1+|L|)^{\deg Q_\lambda}\right).
\]
If the inverse is already a finite expression and no omitted nonzero
block exists, the remainder is instead exactly zero. A program must not
claim a nonzero next term merely because the support semigroup allows one.

\section{The logarithmic example to arbitrary order}
\label{sec:logexample}

Here $a=p=1$, $\delta_1=1$, and $P_1(L)=L+1$. Thus
\begin{equation}\label{eq:logcoeff}
 g_1(y)=y+\sum_{n\ge1}y^{n+1}q_n(\log y),\qquad
 q_n(L)=\frac{(-1)^n}{n!}
 \prod_{j=1}^{n-1}(D_L+n+1+j)(L+1)^n.
\end{equation}
Each $q_n$ has degree $n$. The first five are
\begin{align}
 q_1={}&-L-1,\notag\\
 q_2={}&2L^2+5L+3,\notag\\
 q_3={}&-5L^3-\frac{41}{2}L^2-27L-\frac{23}{2},\notag\\
 q_4={}&14L^4+\frac{241}{3}L^3+\frac{335}{2}L^2+151L+\frac{299}{6},
 \label{eq:logcoefftable}\\
 q_5={}&-42L^5-\frac{3727}{12}L^4-\frac{5365}{6}L^3
 -1256L^2-\frac{5177}{6}L-\frac{2789}{12}.\notag
\end{align}
Further exact coefficients are included in the machine-readable validation
results and can be generated at any chosen cutoff subject to resource
limits.

\subsection{Catalan numbers in the leading logarithms}

The derivative part of $D_L+c$ lowers logarithmic degree, so the leading
coefficient of $q_n$ comes from multiplying only the scalar parts:
\[
 [L^n]q_n(L)
 =\frac{(-1)^n}{n!}\prod_{j=1}^{n-1}(n+1+j)
 =(-1)^n\frac1{n+1}\binom{2n}{n}.
\]
The leading logarithms are signed Catalan numbers. Lower logarithmic
powers are not optional decorations: their derivatives contribute to
higher inverse coefficients. Replacing $\log g$ by $\log y$ too early
recovers some leading terms but misses those subleading contributions.

\subsection{A compact direct Wolfram construction}

The observable homotopy formula immediately gives a minimal coefficient
constructor for the identity-leading case:
\begin{lstlisting}
(* h is the correction to the identity, as an expression in z. *)
identityInverseJet[h_, {z_, y_, n_Integer}] :=
 y + Total[Table[
   (-1)^k/k! D[h^k, {z, k - 1}] /. z -> y,
   {k, 1, n}
 ]];

identityInverseJet[z^2 (1 + Log[z]), {z, y, 4}]
\end{lstlisting}
This finite derivative formula explains why the problem is computationally
tractable. It does not by itself implement mixed-weight sorting, input
normalization, branch metadata, or a rigorous remainder. The full package
adds those layers.

\subsection{The remainder can be seen numerically and symbolically}

For the two-block approximation $A_3=y-y^2(1+\log y)$,
\[
 \frac{g_1(y)-A_3(y)}{y^3\log^2 y}\longrightarrow2.
\]
The independently computed ratios at $y=10^{-20},10^{-50},10^{-100}$ are
approximately $1.892841$, $1.956797$, and $1.978342$, respectively. The slow
approach reflects the next coefficient $5/\log y$, not a failure of the
expansion. The quotient by $y^3$ itself diverges like $2\log^2y$, directly
showing why a bare $\OO(y^3)$ would be incorrect.

\section{The irrational-power example and its convergence radius}
\label{sec:irrexpansion}

Consider the entire family
\[
 f_r(x)=x+x^r,\qquad r>1,\qquad \delta=r-1.
\]
Here the logarithmic polynomial is $1$, so the differential operators in
\eqref{eq:Qformula} act by their constant parts. For $n\ge1$,
\begin{align}\label{eq:fuss}
 c_n
 &=\frac{(-1)^n}{n!}\prod_{j=1}^{n-1}(1+n\delta+j)\notag\\
 &=\frac{(-1)^n}{n}\binom{nr}{n-1}
 =\frac{(-1)^n}{1+n(r-1)}\binom{nr}{n}.
\end{align}
With $c_0=1$, the inverse is
\begin{equation}\label{eq:irrall}
 g_r(y)=\sum_{n=0}^{\infty}c_n y^{1+n(r-1)}.
\end{equation}
There is no rationality requirement on $r$ in this identity.

The first coefficients are
\[
 1,\quad -1,\quad r,\quad-\frac{r(3r-1)}2,
 \quad\frac{r(4r-1)(4r-2)}6,\quad\ldots.
\]
For $r=2$ these reduce to the signed Catalan coefficients of the inverse
of $x+x^2$. For $r=\sqrt2$, equation~\eqref{eq:irranswer} follows at once.

\subsection{A genuine convergent ordinary series in a new parameter}

Put $t=y^{r-1}$ and $v=g_r(y)/y$. The equation becomes
\begin{equation}\label{eq:vpower}
 v+t v^r=1,
\end{equation}
with $v(0)=1$. Thus $v$ is an ordinary analytic function of $t$ near zero,
even though its expression as a function of $y$ uses irrational powers.
The critical point of~\eqref{eq:vpower} satisfies
\[
 v+t v^r=1,\qquad 1+r t v^{r-1}=0.
\]
It occurs at
\[
 v_* =\frac r{r-1},\qquad
 t_*=-\frac{(r-1)^{r-1}}{r^r}.
\]

To establish the radius without assuming that this critical point is the
only relevant complex obstruction, use the exact coefficients. Stirling's
formula gives
\begin{equation}\label{eq:cnstirling}
 |c_n|\sim
 \frac{\sqrt r}{\sqrt{2\pi}(r-1)^{3/2}}\,
 n^{-3/2}\left(\frac{r^r}{(r-1)^{r-1}}\right)^n.
\end{equation}
The root test therefore proves the exact radius
\begin{equation}\label{eq:radius}
 R_t=\frac{(r-1)^{r-1}}{r^r}.
\end{equation}
The corresponding positive-$y$ threshold for this series representation is
$R_y=R_t^{1/(r-1)}$. At $r=\sqrt2$,
\[
 R_t\approx0.4251956066518408,\qquad
 R_y\approx0.1268627051233097.
\]
The real inverse exists for every $y>0$, but this particular Taylor series
in $t$ does not converge for every such $y$. Beyond its convergence disk,
one may continue the analytic solution or solve the monotone real equation
numerically. Global invertibility and the radius of a local representation
are different questions.

\subsection{A non-asymptotic tail majorant inside the disk}

For fixed $r>1$ and any $0<T<R_t$, convergence gives a finite constant
$M_T=\sum_{n\ge0}|c_n|T^n$. If $0\le t<T$, then
\[
 \left|v(t)-\sum_{n=0}^{N-1}c_nt^n\right|
 \le M_T\left(\frac tT\right)^N.
\]
This bound is not optimized, but it separates a uniform convergence
statement from a purely asymptotic order symbol. Sharper alternating-tail
bounds require checking the term-magnitude monotonicity at the particular
$t$; the package does not assume it without a check.

\section{Mixed supports and nonunit leading powers}
\label{sec:mixed}

\subsection{Irrational powers mixed with logarithms}

For
\[
 f(x)=x+x^{\sqrt2}+x^2(1+\log x),
\]
write $r=\sqrt2$ and $L=\log y$. The relative support is generated by
$r-1$ and $1$. The inverse blocks with absolute exponent less than $3$
are shown in Table~\ref{tab:mixed}. They are already sorted by their exact
real exponents.

\begin{table}[htbp]
\centering
\caption{The inverse of $x+x^{\sqrt2}+x^2(1+\log x)$ below power $3$.}
\label{tab:mixed}
\begin{tabular}{@{}ll@{}}
\toprule
Power of $y$ & Coefficient polynomial in $L=\log y$\\
\midrule
$1$ & $1$\\
$r$ & $-1$\\
$2r-1$ & $r$\\
$2$ & $-L-1$\\
$3r-2$ & $-3+r/2$\\
$1+r$ & $(r+2)L+r+3$\\
$4r-3$ & $-4+17r/3$\\
$2r$ & $-(5+3r)L-13/2-5r$\\
\bottomrule
\end{tabular}
\end{table}

The first omitted block is $y^3(2L^2+5L+3)$; a sharpened remainder is
therefore $\OO(y^3L^2)$. The universal bound of
Theorem~\ref{thm:remainder} is more conservative:
$\OO(y^3(1+|L|)^5)$. Both are valid, but they encode different amounts of
support information.

A perturbation-degree cutoff would interleave these terms incorrectly.
For example, repeated use of the $x^{r-1}$ correction can produce several
terms before a single use of a higher-weight correction. The package
therefore enumerates by the inequality
$m_1(r-1)+m_2<B$, not by a single common bound on $m_1+m_2$.

\subsection{Collisions and exact zero blocks}

For $f(x)=x+x^2+x^3$, relative weights $1$ and $2$ are dependent.
The weight-$2$ coefficient receives both a repeated weight-$1$
contribution and a direct weight-$2$ contribution. One obtains
\[
 f^{-1}(y)=y-y^2+y^3+0\,y^4-4y^5+14y^6+\cdots.
\]
The vanishing $y^4$ block is a useful regression test. Keeping one
representative multi-index for each exponent would be wrong: the
contributions must be summed, not deduplicated before coefficient
calculation.

\subsection{A nonunit scale changes both the power and the logarithm}

Consider
\begin{equation}\label{eq:p2example}
 f(x)=4x^2\left(1+x\log x+3x^{\sqrt2}\right).
\end{equation}
Put $u=\sqrt{y/4}$ and $L=\log u=\tfrac12\log(y/4)$. Up to absolute
$u$-power $4$, the inverse is
\begin{align}\label{eq:p2answer}
 g(y)={}&u-\frac12L u^2-\frac32u^{1+\sqrt2}
 +\left(\frac58L^2+\frac14L\right)u^3\notag\\
 &+\frac34\bigl((4+\sqrt2)L+1\bigr)u^{2+\sqrt2}
 +\frac98(3+2\sqrt2)u^{1+2\sqrt2}
 +\OO\!\left(u^4(1+|L|)^3\right).
\end{align}
Replacing $L$ by $\log y$ would change coefficients and make this
expansion false. The leading coefficient $4$ contributes a logarithmic
shift as well as a multiplicative scale. This is why the package carries
$u$ and $L$ as independent internal symbols until final specialization.

\section{Residual checking and the passage to actual inverse errors}
\label{sec:residual}

\subsection{The exact residual identity}

Let $A(y)>0$ be any finite approximation with $A/u\to1$. Define
\[
 R(y)=f(A(y))-y,\qquad
 \rho(u)=\frac{f(A(y))}{a u^p}-1.
\]
These are related by $R=a u^p\rho$. The mean-value theorem gives, for a
point $\xi$ between $A$ and $g$,
\begin{equation}\label{eq:residualtransfer}
 g(y)-A(y)=-\frac{R(y)}{f'(\xi)}.
\end{equation}
Since $A/u\to1$, $g/u\to1$, and $f'(x)\sim apx^{p-1}$,
\begin{equation}\label{eq:relative-transfer}
 g(y)-A(y)=-\frac up\,\rho(u)\,(1+o(1))
\end{equation}
where the ratio is interpreted along nonzero residual values. In
particular, if
\[
 \rho(u)=u^\lambda S(L)+o\bigl(u^\lambda(1+|L|)^{\deg S}\bigr),
 \qquad S\ne0,
\]
then the first inverse correction is $-u^{1+\lambda}S(L)/p$.
The divisor is $p$, not $1$, unless the leading exponent is one.

This proves that a power--log residual calculation can certify the
asymptotic correctness of an approximation when combined with the stated
real-branch and derivative hypotheses. A list of formal zero coefficients
alone is not a finite-precision interval certificate.

\begin{corollary}[Global error bounds for the original examples]
For every $y>0$ and every positive approximation $A$, one has
\begin{align*}
 |g_1(y)-A|&\le
 \frac{|A+A^2(1+\log A)-y|}{1-2e^{-5/2}},\\
 |g_2(y)-A|&\le |A+A^{\sqrt2}-y|.
\end{align*}
\end{corollary}
\begin{proof}
Apply the mean-value theorem and the global bounds
$f_1'\ge1-2e^{-5/2}$ and $f_2'\ge1$ proved in
Section~\ref{sec:question}. No small-argument restriction or root bracket
is needed for these two inequalities, beyond $A,y>0$.
\end{proof}

\subsection{The independent sparse-ring checker}

The package stores $A/u=1+V$ as a sparse set of positive-weight blocks.
It independently constructs
\begin{equation}\label{eq:checker}
 \rho=(1+V)^p\left[1+\sum_i u^{\delta_i}(1+V)^{\delta_i}
 P_i\bigl(L+\log(1+V)\bigr)\right]-1.
\end{equation}
Binomial and logarithmic series are evaluated in a clipped sparse ring,
using~\eqref{eq:unit-calculus}. Polynomial substitution uses Horner's rule.
Each multiplication discards exponents at or beyond the requested relative
cutoff. No coefficient formula from~\eqref{eq:Qformula} is reused in this
checker.

For an inverse generated with absolute cutoff $C$, the expected residual
vanishes below relative power $C-1$. The checker reports a Boolean
certificate only after the requested residual range covers that whole
interval. A shorter range returns an explicit
\code{Missing["InsufficientCutoff"]} rather than an unjustified success.
Increasing the residual cutoff can expose the first omitted inverse block
via~\eqref{eq:relative-transfer}.

\subsection{An explicit finite-point slope bound}

There is also a concrete route from a residual to a numerical error bound.
For $0<u<1/2$, consider $I_u=[u/2,2u]$ and put
$M=|\log u|+\log2$. Define coefficient majorants
\[
 \widehat P_i(M)=\sum_k |[L^k]P_i|M^k,
\]
and, with $T_i=(p+\delta_i)P_i+P_i'$, define $\widehat T_i$ similarly.
Set
\begin{align*}
 E_0(u)&=\sum_i(2u)^{\delta_i}\widehat P_i(M),\\
 E_1(u)&=\sum_i(2u)^{\delta_i}\widehat T_i(M).
\end{align*}
If
\[
 E_0(u)<1-2^{-p},\qquad E_1(u)\le p/2,
\]
then $f(u/2)<a u^p<f(2u)$ and
\begin{equation}\label{eq:explicit-slope}
 f'(x)\ge m_u:=\frac{ap}{2}\,u^{p-1}2^{-|p-1|}
 \quad\text{for }x\in I_u.
\end{equation}
Indeed, the first inequality bounds $1+H(x)$ on the bracket; the second
bounds the correction to $p$ in~\eqref{eq:derivativeasym}. The minimum of
$x^{p-1}$ on this bracket is $u^{p-1}2^{-|p-1|}$.

Once $A\in I_u$, equation~\eqref{eq:residualtransfer} gives
\begin{equation}\label{eq:finitebound}
 |g(y)-A(y)|\le\frac{|f(A(y))-y|}{m_u}.
\end{equation}
To turn this into a machine-certified numerical bound, all inequalities
and function evaluations must themselves use verified enclosures. The
package's algebraic residual checker does not silently claim that an
ordinary floating-point comparison supplies such an enclosure.

\subsection{Newton refinement as a scalable alternative}

In the normalized ring, $\partial\mathcal F/\partial V$ has constant term
$p$ and is a unit. A Newton step
\[
 V_{\mathrm{new}}=V-
 \frac{\mathcal F(V)}{\partial\mathcal F(V)/\partial V}
\]
therefore makes sense with clipped series arithmetic. If the error has
positive valuation $\lambda$, the next error has valuation at least
$2\lambda$, because the normalized equation is analytic in $V$ and its
linearization is invertible. This gives a Newton--Hensel strategy with
increasing cutoffs.

The shipped generator uses the explicit multi-index formula; its
independent checker already supplies most of the sparse operations needed
for a future Newton backend. Newton generation is described here as an
alternative algorithm, not claimed as an additional implemented public API.

\section{Input jets, observables, and information below all powers}
\label{sec:extensions}

\subsection{When the finite input is an approximation}

Suppose $F$ is a prepared finite germ and an actual function satisfies
\begin{equation}\label{eq:inputjet}
 f(x)-F(x)=\OO\!\left(x^{p+\beta}(1+|\log x|)^k\right),
 \qquad \beta>0,
\end{equation}
with both functions having the addressed inverse branches and
$f'(x)\sim apx^{p-1}$. Evaluating $f$ at $F^{-1}(y)$ and using the residual
transfer theorem yields
\begin{equation}\label{eq:inputjetout}
 f^{-1}(y)-F^{-1}(y)
 =\OO\!\left(u^{1+\beta}(1+|\log u|)^k\right).
\end{equation}
Thus inverse terms above that uncertainty scale are not determined by the
supplied input jet alone. The package interprets its finite expression as
exact; a user applying it to a truncated jet must combine its output
remainder with~\eqref{eq:inputjetout}.

A value asymptotic does not automatically imply a derivative asymptotic.
For example,
\[
 f(x)=x+x^4\sin(x^{-4})=x+\OO(x^4)
\]
has
\[
 f'(x)=1+4x^3\sin(x^{-4})-4x^{-1}\cos(x^{-4}),
\]
which is not eventually positive. One cannot infer an inverse branch
merely by differentiating the written order symbol. Exact finite prepared
input avoids this problem because its derivatives are explicit.

\subsection{Powers and logarithms of the inverse}

Replace the observable $\Phi(Z)=Z^{1/p}$ in~\eqref{eq:LB} by $Z^{q/p}$.
For fixed real $q$, the same calculation gives
\begin{equation}\label{eq:observable}
 g(y)^q=u^q\left(1+\sum_{\boldsymbol m\ne0}u^\Delta
 \frac{q(-1)^n}{p^n\boldsymbol m!}
 \prod_{j=1}^{n-1}(D_L+q+\Delta+pj)P^{\boldsymbol m}(L)\right).
\end{equation}
For $q=0$ the correction vanishes and the identity is $1=1$.
Differentiating this identity with respect to $q$ at zero gives
\begin{equation}\label{eq:loginverse}
 \log g(y)=L+\sum_{\boldsymbol m\ne0}u^\Delta
 \frac{(-1)^n}{p^n\boldsymbol m!}
 \prod_{j=1}^{n-1}(D_L+\Delta+pj)P^{\boldsymbol m}(L).
\end{equation}
These are useful when the quantity of interest is an observable of the
inverse rather than the inverse itself. They also provide independent
symbolic consistency tests.

For a stored inverse block $u^eQ(L)$,
\[
 \frac{\dd}{\dd y}\bigl(u^eQ(L)\bigr)
 =\frac1{ap}u^{e-p}(e+D_L)Q(L).
\]
Termwise differentiation is justified within the analytic finite-input
construction on an appropriate local sector; it is not licensed merely by
a bare real-axis big-$\OO$ remainder for an arbitrary function.

\subsection{Flat perturbations have separate inverse effects}

Let $\widetilde f=f+\varepsilon$ where $\varepsilon$ is smaller than every
power of $x$, and suppose the perturbed function retains the addressed
branch and the required derivative control. Then
\[
 \widetilde g(y)-g(y)
 =-\frac{\varepsilon(g(y))}{f'(g(y))}(1+o(1))
\]
under the corresponding stability hypotheses, for example when the
perturbation and its derivative are uniformly negligible on the small
comparison interval. Such an effect is invisible to every finite
power--log cutoff. This is not a failure of reversion; it is information
not present in the chosen formal scale.

For instance, $x$ and $x+e^{-1/x}$ have the same pure-power asymptotic
expansion at zero, but their inverses differ by
$-e^{-1/y}(1+o(1))$. The perturbation is monotone and its derivative is flat,
so the real-germ comparison is well behaved. An exponential sector would
be needed to record this difference explicitly.

\section{Leading logarithmic pivots and other endpoints}
\label{sec:leadinglog}

\subsection{The exact inverse of a leading log-power monomial}

The core class excludes a nonconstant logarithmic \emph{leading} block.
That case should be normalized differently, not forced into a correction
of positive power valuation. Consider
\begin{equation}\label{eq:leadinglog}
 f_0(x)=a x^p(-\log x)^q,\qquad 0<x<1,\quad a,p>0.
\end{equation}
With $t=-\log x$ and $Y=y/a$, the equation is
\[
 Y=t^q e^{-pt}.
\]
For $q\ne0$, this gives
\begin{equation}\label{eq:Wbase}
 b(y)=\exp\left[\frac qp
 W_k\left(-\frac pq\left(\frac ya\right)^{1/q}\right)\right].
\end{equation}
The branch is forced by $t\to+\infty$:
\[
 k=-1\quad(q>0),\qquad k=0\quad(q<0).
\]
For $q=0$, use $b(y)=(y/a)^{1/p}$. The real branches of Lambert $W$ and
their domains are described in DLMF~\S4.13~\cite{dlmfw}.

For $q>0$, the small-$x$ branch has
$0<Y<(q/(pe))^q$. The principal branch $W_0$ would instead give $t\to0$
and $x\to1$, another genuine real inverse germ of the same leading model.
For $q<0$, the argument of $W_0$ is positive and tends to infinity, yielding
the desired $t\to\infty$ branch. This branch distinction is precisely why
\code{LogPowerBaseInverse} is a separate helper.

\subsection{The logarithm-of-logarithm scale}

Put $H=\log(a/y)$. The equation $pt=H+q\log t$ gives, by successive
substitution,
\begin{equation}\label{eq:tlogexp}
 t=\frac1p\left[
 H+q\log\frac Hp+
 \frac{q^2}{H}\log\frac Hp+
 \OO\!\left(\frac{(\log H)^2}{H^2}\right)\right].
\end{equation}
Consequently,
\begin{equation}\label{eq:xlogexp}
 b(y)=\left(\frac ya\right)^{1/p}
 \left(\frac Hp\right)^{-q/p}
 \left[1-\frac{q^2}{pH}\log\frac Hp+
 \OO\!\left(\frac{(\log H)^2}{H^2}\right)\right].
\end{equation}
The relative small quantities are now inverse logarithms with
logarithm-of-logarithm coefficients. These are not positive powers of the
uniformizer in Definition~\ref{def:class}. The formula identifies the
necessary extension of scale instead of pretending that the existing
finite polynomial-log representation already covers it.

\subsection{Perturbing an exactly invertible pivot}

Suppose $f=f_0+h$ and the addressed inverse $b=f_0^{-1}$ is known. In the
coordinate $s=f_0(x)$ the equation becomes
\[
 y=s+\eta(s),\qquad \eta(s)=h(b(s)),\qquad x=b(s).
\]
The observable formula gives the homotopy jet
\begin{equation}\label{eq:pivotjet}
 x=b(y)+\sum_{n=1}^{N}\frac{(-1)^n}{n!}
 \frac{\dd^{n-1}}{\dd y^{n-1}}\left(b'(y)\eta(y)^n\right).
\end{equation}
This is implemented by \code{InversePerturbationJet}. The argument supplied
as the perturbation is $\eta$ in the \emph{pivot coordinate}, not the
original $h(x)$. Establishing the asymptotic ordering, convergence, or
remainder of such a jet requires the relevant prepared-scale hypotheses;
the helper deliberately does not attach an automatic big-$\OO$ claim.

\subsection{Shifts, the other real side, and infinity}

For a finite endpoint $x_0$ with $f(x_0)=y_0$, first introduce
$x=x_0+t$ and $y=y_0+s$ and address the intended side $t\to0^+$ or
$t\to0^-$. Real noninteger powers on a negative side require an explicit
real-valued definition; a complex principal power must not be silently
reinterpreted as a real one.

For an inverse at infinity, a useful transformation is
\[
 F(t)=\frac1{f(1/t)}.
\]
If $G=F^{-1}$ is the addressed inverse near zero, then
\[
 f^{-1}(y)=\frac1{G(1/y)}.
\]
The transformed function may have an infinite or rational power--log
expansion rather than exact finite input. In that case one must prepare a
finite jet and propagate its uncertainty as in
Section~\ref{sec:extensions}. These endpoint transformations are
mathematical recipes, not automatic preprocessing implemented by the core
parser.

\section{Wolfram Language package: interface and implementation}
\label{sec:package}

\subsection{Loading and solving the original examples}

After extracting the archive, load the package from its actual path:
\begin{lstlisting}
Get["/path/to/PowerLogInverse/Kernel/PowerLogInverse.wl"];

PowerLogInverse[
 x + x^2 (1 + Log[x]), {x, y, 5}
]

PowerLogInverse[
 x + x^Sqrt[2], {x, y, 4 Sqrt[2] - 3}
]
\end{lstlisting}
These requests mean strict absolute uniformizer-power cutoffs. Since
$a=p=1$ in both examples, the uniformizer is just $y$. The second call
retains exactly the four blocks in~\eqref{eq:irranswer}.

\noindent\begin{minipage}{\linewidth}
For the coefficients and error metadata together:
\begin{lstlisting}
d = PowerLogInverseData[
 x + x^2 (1 + Log[x]), {x, y, 5}
];
d["Expression"]
d["Remainder"]
(* PowerLogOrder[y, 5, 4] *)

PowerLogResidual[d]["Certificate"]
(* Expected in a native kernel: True *)

PowerLogResidual[d, 5]["Expression"]
(* First nonzero relative residual blocks, below u-power 5. *)
\end{lstlisting}
\end{minipage}

The finite expression is the approximation only. It is not an expression
with an order object added to it. The descriptor
\code{PowerLogOrder[u,C,K]} means
$\OO(u^C(1+|\log u|)^K)$ as $u\to0^+$. No algebraic rules for adding,
multiplying, or composing such descriptors are installed.

\subsection{Public functions}

\begin{longtable}{@{}p{0.36\textwidth}p{0.60\textwidth}@{}}
\toprule
Function & Contract\\
\midrule
\endhead
\code{PowerLogInverse} & Parse and normalize a finite input expression;
return its finite positive-real inverse approximation.\\[4pt]
\code{PowerLogInverseData} & Return the same approximation with exact
support blocks, normalization rules, remainder, input conditions, and
resource metadata.\\[4pt]
\code{PowerLogReversion} & Low-level prepared-input interface
$[a,p,\{\{\delta_i,P_i(L)\}\},\{u,L\},C]$; returns structured data.\\[4pt]
\code{PowerLogResidual} & Independently expand the normalized residual of
the stored finite approximant through a relative cutoff.\\[4pt]
\code{PowerLogOrder} & Inert meaning-bearing remainder descriptor, not a
series arithmetic head.\\[4pt]
\code{LogPowerBaseInverse} & Return the exact leading-log pivot inverse
with the selected real \code{ProductLog} branch.\\[4pt]
\code{InversePerturbationJet} & Return a finite observable Lagrange homotopy
jet for a supplied pivot coordinate; no automatic remainder.\\
\bottomrule
\end{longtable}

The low-level interface is useful when normalization is already known:
\begin{lstlisting}
d = PowerLogReversion[
 4, 2, {{1, L}, {Sqrt[2], 3}}, {u, L}, 4
];
d["Terms"]
PowerLogResidual[d]["Certificate"]
\end{lstlisting}
Here \code{d["Expression"]} still contains the independent symbols $u,L$.
Specialize them to $u=\sqrt{y/4}$ and $L=\tfrac12\log(y/4)$ only when an
expression in $y$ is desired. The high-level interface performs that
specialization and retains \code{UniformizerRule} and \code{LogRule}.

\subsection{The parser's deliberate domain}

The input must expand into a finite sum of $x^r$ times polynomials in
\code{Log[x]}. Exponents, including the leading $p$ and the cutoff, must
be exact real algebraic numbers. Coefficients must be exact and provably
real constants. The smallest power must have a constant positive
coefficient, and its exponent must be positive.

The parser rejects a nonconstant leading logarithmic block, inverse
logarithmic powers, unprepared functions such as $\sin x$, unresolved
symbolic exponents, and machine-precision exponents. These restrictions
are explicit interface boundaries, not claims that the corresponding
mathematical problems lack asymptotic expansions. A symbolic parameter can
be specialized before calling the package, or one can use the general
homotopy helper when its hypotheses have been established separately.

A \code{ConditionalExpression} is accepted only when its condition is
provably implied by $x>0$. An unresolved condition is returned as a
\code{Failure}; it is not discarded. Eventual conditions such as $x<1$
that are not globally implied by $x>0$ require explicit germ preparation.
Use distinct source and target symbols, with the target absent from the
input expression.

No use is made of \code{PowerExpand}. Identities involving logarithms and
powers are derived from positivity and the stated normalization, not from
unconditional manipulation of principal complex branches. Likewise, no
irrational exponent is replaced by a rational approximation. If
$r$ is changed to $r+\epsilon$, the ratio $x^{r+\epsilon}/x^r=
\exp(\epsilon\log x)$ need not remain close to $1$ in the asymptotic limit,
no matter how small a fixed nonzero $\epsilon$ is.

\subsection{Exact exponent keys and collisions}

Exponent arithmetic is reduced with \code{RootReduce}. Rows are sorted by
exact inequalities and adjacent equal exponents are combined using exact
equality tests. This makes collision handling independent of accidental
syntactic differences between equivalent algebraic numbers. An undecided
order is a failure, not an invitation to guess from decimal approximations.
Logarithmic polynomial coefficients are expanded and simplified before
zero blocks are removed.

Internally, a row is \code{\{exponent, polynomial\}}. Lists rather than
symbolically keyed associations are used for the sparse arithmetic.
Associations provide the outer API metadata. The residual checker uses
the stored \code{Terms} and prepared corrections; it does not infer new
coefficients from the already computed inverse formula.

\subsection{Enumeration and resource controls}

For the requested relative cutoff $B$, the generator enumerates only the
integer simplex $\boldsymbol m\cdot\boldsymbol\delta<B$. Generators whose
weight is already at least $B$ can be omitted from coefficient enumeration,
though they remain part of the input model and its remainder analysis.
The implementation uses an iterative depth-first traversal, avoiding a
dependence on the Wolfram kernel's recursion-depth setting.

The public option names are strings. Their defaults are
\begin{center}
\begin{tabular}{@{}lr@{}}
\toprule
Option & Default\\
\midrule
\code{"MaxMultiIndices"} & $100000$\\
\code{"MaxTerms"} & $50000$\\
\code{"MaxTotalDegree"} & $256$\\
\code{"MaxProducts"} & $2000000$\\
\bottomrule
\end{tabular}
\end{center}
The product limit also bounds a growing enumeration frontier. Exceeding a
budget returns a named \code{Failure}, not a partially computed expression
masquerading as a complete truncation. These limits do not constitute
wall-clock guarantees: exact algebraic simplification and coefficient
expression growth can still dominate a calculation.

For fixed $s$ and positive generators, the number of visited terminal
indices is bounded by a finite box and grows like the volume of a simplex,
$B^s/(s!\prod_i\delta_i)$, at large cutoffs. When generators are dependent,
the number of distinct output exponents can be much smaller than the
number of contributing multi-indices. Direct coefficient enumeration is
transparent and useful at moderate cutoffs; a sparse Newton backend is
an appropriate optimization for large, highly colliding problems.

\subsection{Implementation map}

The core file separates the mathematical and operational layers:
\begin{enumerate}[leftmargin=1.8em,itemsep=0.2em]
\item Exact-domain validation, polynomial parsing, positive leading-term
selection, and condition checking.
\item Cutoff-finite multi-index enumeration and the differential-polynomial
coefficient formula.
\item Exact exponent collection, finite-expression construction, and the
analytic remainder metadata from Theorem~\ref{thm:remainder}.
\item Independent clipped sparse addition, multiplication, binomial powers,
logarithms, polynomial substitution, and residual checking.
\end{enumerate}
The last two public helpers are intentionally outside the finite core:
they expose leading-log pivots and observable homotopy derivatives without
pretending that their general remainder problem has already been solved
by the core parser.

\section{Verification evidence and reproducibility}
\label{sec:validation}

\subsection{What was executed}

The archive contains an independent Python implementation using exact
SymPy algebra and high-precision mpmath arithmetic. It passed
\textbf{46 of 46 checks}. The checks cover the explicit logarithmic
coefficients, generalized Fuss--Catalan coefficients, ordinary Catalan
reversion, strict irrational cutoff boundaries, mixed supports, dependent
and duplicate generators, nonunit and fractional leading exponents,
exact residual cancellation, deliberate coefficient corruption, and
numerical improvement under successive truncation.

The Python coefficient generator uses a bounded Cartesian enumeration,
whereas the Wolfram implementation uses a simplex depth-first traversal.
The Python residual checker uses clipped binomial and logarithmic
composition rather than the coefficient formula. An additional ordinary
integer-power substitution checks the logarithmic example independently
of that sparse residual code. These checks test the mathematics by several
routes; they do not amount to execution of the Wolfram source.

The archive also contains \textbf{40 native Wolfram Language regression
test definitions}. Native Wolfram-kernel execution was not available in
the build environment. The Wolfram files passed a lexical audit for
balanced delimiters, strings, and nested comments, but that audit is not a
Wolfram parser or evaluator. Accordingly, native execution is recorded as
\emph{unverified}, not as a successful test run.

\subsection{Numerical results}

Table~\ref{tab:numerical} reports absolute inverse errors at $y=10^{-8}$.
Roots were obtained independently by high-precision bisection in a real
bracket. These are numerical diagnostics, not interval-certified error
bounds. The exact coefficients were converted to the requested precision
only at evaluation time.

\begin{table}[htbp]
\centering
\caption{Independently measured errors at $y=10^{-8}$. Cutoffs refer to $u$.}
\label{tab:numerical}
\begin{tabular}{@{}lll@{}}
\toprule
Input function & Strict cutoff $C$ & Absolute inverse error\\
\midrule
$x+x^2(1+\log x)$ & $3$ & $5.8953980\times10^{-22}$\\
 & $5$ & $1.1639260\times10^{-34}$\\
 & $7$ & $3.0778899\times10^{-47}$\\[2pt]
$x+x^{\sqrt2}$ & $1+2(\sqrt2-1)$ & $3.3323930\times10^{-15}$\\
 & $1+4(\sqrt2-1)$ & $2.2301916\times10^{-21}$\\
 & $1+6(\sqrt2-1)$ & $1.8441485\times10^{-27}$\\[2pt]
Equation~\eqref{eq:p2example} & $3$ & $7.2770394\times10^{-12}$\\
 & $4$ & $5.4503639\times10^{-15}$\\
 & $5$ & $4.5863230\times10^{-18}$\\
\bottomrule
\end{tabular}
\end{table}

The first two-block irrational approximation is relatively less accurate
than the two-block logarithmic approximation at the same $y$. This is
expected: its first relative small scale is $y^{\sqrt2-1}$, which decreases
much more slowly than $y|\log y|$. Counting blocks alone is not a fair
measure of asymptotic accuracy across different supports.

\subsection{Reproducing the computations}

From the archive root, the independent validation commands are
\begin{lstlisting}[language={}]
python Validation/validate.py
python Validation/static_check.py
\end{lstlisting}
The former regenerates \code{validation-results.json}, including numerical
errors and coefficients; the latter regenerates
\code{static-results.json}. The Python requirements are SymPy and mpmath.
Versions used in the executed run are recorded in the JSON file.

For a native Wolfram installation:
\begin{lstlisting}[language={}]
wolframscript -file Tests/run-tests.wls
wolframscript -file Examples/examples.wls
\end{lstlisting}
Alternatively, evaluate \code{TestReport} on the supplied \code{.wlt}
file in a notebook. The test suite includes failure-path tests in addition
to coefficient comparisons. In particular, a corrupted inverse block must
make the independent residual certificate false, and an insufficient
residual range must not return true.

Compile this self-contained article source with
\begin{lstlisting}[language={}]
latexmk -pdf -interaction=nonstopmode -halt-on-error \
  real_branch_reversion.tex
\end{lstlisting}
The archive includes the compiled PDF, so no TeX installation is needed
to read the article.

\section{Scope and conclusions}
\label{sec:conclusion}

The two motivating inverses are covered to arbitrary finite order by the
same coefficient calculus. Their apparent incompatibility with ordinary
series reversion comes from representation and branch selection, not
from a lack of constructive formulas. The first needs logarithmic
polynomials at integer powers; the second needs irrational powers; their
mixtures need a finitely generated positive support semigroup.

The finite-input result is stronger than a formal expansion: sufficiently
near the selected positive endpoint, the constructed power--log series
actually converges. Its strict cutoff has a proved, computable logarithmic
remainder envelope. An independently computed normalized residual
provides another route to asymptotic verification and, with a verified
slope bound, to finite-point error control.

The current executable core is deliberately exact and finite. It does not
automatically discover arbitrary transseries scales, decide symbolic
parameter orderings, invert every leading logarithmic polynomial, classify
all global branches, or certify floating-point intervals. The article
identifies concrete extensions rather than hiding those tasks inside an
unqualified universal inverse command. Leading logarithmic monomials have
an implemented real-$W$ pivot helper; general pivot perturbations have an
implemented homotopy-jet helper; the broader scale and residual hypotheses
remain explicit.

The essential computational contract is simple: specify the real germ,
prepare the leading term, enumerate the actual ordered support, collect
all equal-weight contributions, and attach only a remainder justified by
that preparation. This is sufficient to resolve both original examples
and a substantial mixed power--logarithmic class without invoking an
unspecified inverse branch.

\appendix
\section{A coefficient-level pseudocode specification}

The following pseudocode describes the generator independently of Wolfram
Language evaluation details. It is also a compact reference for a port to
another exact symbolic system.
\begin{lstlisting}[language={}]
Input: a > 0, p > 0, corrections (delta_i, P_i), cutoff C > 1
Normalize duplicate deltas; remove zero P_i
u := positive (y/a)^(1/p); L := Log(y/a)/p
B := C - 1
rows := [(1, 1)]

for each nonzero multi-index m with dot(m, delta) < B:
    n := sum(m)
    Delta := dot(m, delta)
    Q := product(P_i(L)^m_i)
    for j := 1 to n-1:
        Q := derivative(Q, L) + (1 + Delta + p*j)*Q
    Q := (-1)^n * Q / (p^n * product(factorial(m_i)))
    append (1 + Delta, Q) to rows

Sort rows by exact real exponent
Add every polynomial sharing the same exponent
Remove exact zero blocks
Return sum(u^e * Q(L) for (e,Q) in rows)

If corrections are nonempty:
    N := ceiling(B / min(delta_i))
    K := N * max(degree(P_i))
    remainder := O(u^C * (1 + abs(Log(u)))^K)
Else:
    remainder := 0
\end{lstlisting}

An implementation must test both domain predicates and cutoff comparisons
before relying on this specification. Algebraic real exponents supply an
exact decision domain. Numerical approximations of these comparisons
change the specification and are not used by the shipped package.

\section{The first logarithmic coefficient recurrences as a check}

For the identity-leading logarithmic example, write
$g=y+\sum_{n\ge1}q_n(L)y^{n+1}$. Let $V=g/y-1$. Then
\[
 0=V+y(1+V)^2\bigl(1+L+\log(1+V)\bigr).
\]
The coefficient of $y$ gives $q_1=-(1+L)$. At the next relative power,
\[
 q_2+(2L+3)q_1=0,
\]
so $q_2=(2L+3)(L+1)=2L^2+5L+3$. One order further,
\[
 q_3+(2L+3)q_2+\left(L+\frac52\right)q_1^2=0,
\]
which gives the $q_3$ displayed in~\eqref{eq:logcoefftable}. These equations
come directly from substitution and provide a transparent low-order
check on the differentiated Lagrange formula.

\section{Archive contents and provenance}

\begin{longtable}{@{}p{0.49\textwidth}p{0.47\textwidth}@{}}
\toprule
Path & Purpose\\
\midrule
\endhead
\path{article/real_branch_reversion.tex} & Self-contained article source.\\
\path{article/real_branch_reversion.pdf} & Compiled article.\\
\path{Kernel/PowerLogInverse.wl} & Main Wolfram Language implementation.\\
\path{Kernel/init.m} & Standard package loader.\\
\path{Tests/PowerLogInverse.wlt} & Native Wolfram regression definitions.\\
\path{Tests/run-tests.wls} & Native test runner with an exit status.\\
\path{Examples/examples.wls} & Worked package calls and a numerical diagnostic.\\
\path{Validation/validate.py} & Executed independent exact/numerical checks.\\
\path{Validation/validation-results.json} & Reproducible results, coefficients, and versions.\\
\path{Validation/static_check.py} & Lexical delimiter audit, not a native parser.\\
\path{Validation/static-results.json} & Audit results.\\
\path{Sources/question.md} & Text of the user-supplied archived question.\\
\path{README.md}, \path{VALIDATION.md} & Usage and verification-status notes.\\
\path{LICENSE} & MIT No Attribution license for the newly prepared material.\\
\bottomrule
\end{longtable}

The mathematical source of the problem is the user-supplied archive of
Mathematica Stack Exchange question 236367. The ProveIt corpus was consulted
for its formal/analytic architecture, particularly the inversion volume's
attention to support, branches, and residual estimates. Its main source
file was identified by Git blob
\texttt{aec3c70d4793e286b462e7d8894dede7c754995e} during this work. No claim
is made that the entire evolving corpus was audited or that the package
is a formal proof-assistant implementation of it.

\clearpage
\begin{thebibliography}{9}

\bibitem{question}
Vladimir Reshetnikov,
\emph{How to get an asymptotic of the real-valued branch of the inverse function?}
Mathematica Stack Exchange, question 236367.
User-supplied archive consulted on 7 September 2026.
\url{https://mathematica.stackexchange.com/q/236367}.

\bibitem{proveit}
Vladimir Reshetnikov and contributors,
\emph{ProveIt: series-and-transseries draft corpus}, particularly
\path{Transseries_And_Inversion/transseries_and_inversion.tex}.
Working repository, consulted 7 September 2026.
\url{https://github.com/VladimirReshetnikov/ProveIt/tree/main/Analysis/FabiusFunction/docs/semi-formalized-research-frontiers/drafts/series-and-transseries}.

\bibitem{dlmfinverse}
NIST Digital Library of Mathematical Functions,
\S1.10(vii), \emph{Inverse Functions}, including the Lagrange inversion
and extended inversion theorems.
\url{https://dlmf.nist.gov/1.10#vii}.

\bibitem{dlmfw}
NIST Digital Library of Mathematical Functions,
\S4.13, \emph{Lambert $W$-Function}.
\url{https://dlmf.nist.gov/4.13}.

\bibitem{vdh}
Joris van der Hoeven,
\emph{Transseries and Real Differential Algebra},
Lecture Notes in Mathematics 1888, Springer, 2006.
\url{https://doi.org/10.1007/3-540-35590-1}.

\bibitem{wlseries}
Wolfram Research,
\emph{SeriesData} and \emph{InverseSeries}, Wolfram Language documentation,
consulted 7 September 2026.
\url{https://reference.wolfram.com/language/ref/SeriesData.html};
\url{https://reference.wolfram.com/language/ref/InverseSeries.html}.

\bibitem{wlinverse}
Wolfram Research,
\emph{InverseFunction}, Wolfram Language documentation,
consulted 7 September 2026.
\url{https://reference.wolfram.com/language/ref/InverseFunction.html}.

\end{thebibliography}
\end{document}
